\section{Sparse equations over fixed fields and Frobenius multiplicities}
\label{sec:ff}
This section separates three ways in which a sparse equation can conceal a counting problem: a varying coefficient field, a fixed coefficient field with roots in its algebraic closure, and a smooth curve specified by its function field. The distinction is essential because the size of the relevant cohomology can be exponential even when the equation has polynomial bit length.

For a ground field $k$, the subscript $\lambda$ in $\mathbb Q_\lambda$ or $\overline{\mathbb Q}_\lambda$ denotes a prime invertible in $k$. Cohomology with these coefficients is geometric \'etale cohomology, after base change to an algebraic closure of $k$; the subscript $c$ denotes compact supports and $(r)$ denotes the $r$th Tate twist. Coherent cohomology is written with its $\mathcal O$ coefficients.

The complexity of point counting already depends on these representation choices in earlier work. Von zur Gathen, Karpinski, and Shparlinski (1996, Theorem~4.7, Section~4; author's version p.~15) prove $\#\mathrm P$-completeness for sparse affine plane-curve counting over an input finite field under probabilistic polynomial-time parsimonious Turing reductions \cite{ff-gathen}. Lauder and Wan (2008, Theorem~1; preprint p.~2) compute the zeta function of an affine hypersurface of degree $d$ in $n$ variables over $\mathbb F_{p^a}$ in $P(p^nd^{n^2}a^n)$ bit operations for an explicit polynomial $P$ \cite{ff-lauder-wan}. The bound is polynomial in the degree and extension degree when $n,p$ are fixed; it need not be polynomial in a binary exponent or a circuit encoding. Harvey (2015, Theorem~1.1) fixes an arithmetic scheme $X$ and gives deterministic bounds $p\log^{1+o(1)}p$ or $p^{1/2}\log^{2+o(1)}p$ for one reduction, and $R\log^{3+o(1)}R$ for all primes below $R$ \cite{ff-harvey}. His constants depend on $X$; Theorems~1.2--1.3 display the dependence on dimension and degree for hypersurfaces. These algorithms and the reductions below concern compatible, but different, ranges of input parameters.

Milovanov (2019, Theorem~1, p.~77, and Section~3, pp.~78--79) proves a deterministic parsimonious reduction from counting satisfying assignments to counting roots of a sparse univariate polynomial over $\mathbb F_{p^b}$, for every fixed prime $p$ \cite{ff-milovanov}. Here $b$ is written in unary, coefficients are elements of that extension field, and exponents are binary integers. His earlier preprint treats characteristic two. Cao, Feng, Gao, Huang, and Qiu (2026, Lemma~7 and Proposition~9, pp.~4--7) instead fix the coefficient field and seek common roots in its algebraic closure \cite{ff-cao}. Their separating auxiliary root of unity packs many Boolean tests into one sparse polynomial. The following construction uses that packing and gives every satisfying assignment the same number of geometric representatives.

\subsection{Balanced roots and a sparse repeated-root locus}
Let $\Phi$ be a width-three CNF with $n\geq1$ declared variables and $m\geq1$ clauses, and let $S$ be its number of satisfying assignments. Appending a new variable with a positive unit clause enforces these harmless size conventions parsimoniously. Repeated literals may first be simplified, and an empty clause is allowed. Put $M=m+n$.

Choose distinct primes $\ell_1,\ldots,\ell_n\equiv1\pmod8$, and put $D=\prod_i\ell_i$. For the nonzero squares $R_i$ modulo $\ell_i$ define, over $\mathbb F_2$,
\[
 C_i(U)=\sum_{a\in R_i}U^a,
 \qquad \Gamma_i(U)=\sum_{a=0}^{\ell_i-1}U^a.
\]
On $\mu_{\ell_i}$, the group of $\ell_i$th roots of unity, $\Gamma_i$ is one at $1$ and zero elsewhere. Because $2$ is a square modulo $\ell_i$, squaring permutes $R_i$, so $C_i(u)^2=C_i(u)$ for every $u\in\mu_{\ell_i}$. Thus $C_i$ takes Boolean values. If $b$ is a nonsquare modulo $\ell_i$, then for $u\neq1$,
\[
 C_i(u)+C_i(u^b)=\sum_{a=1}^{\ell_i-1}u^a=1.
\]
Raising to $b$ exchanges the two fibers. Each fiber on the primitive $\ell_i$th roots therefore has size $(\ell_i-1)/2$.

The homomorphism
\[
 \mu_D\longrightarrow\prod_i\mu_{\ell_i},
 \qquad y\longmapsto (y^{D/\ell_i})_i
\]
is an isomorphism: its kernel is trivial, since the exponents $D/\ell_i$ have gcd one, and both groups have $D$ elements. The simultaneous guards $\Gamma_i(y^{D/\ell_i})=0$ select the primitive $D$th roots. Among them every Boolean word has exactly
$\prod_i(\ell_i-1)/2$ representatives.

For each clause take the product of its violation factors: a positive literal $x_i$ contributes $1+C_i(Y^{D/\ell_i})$, and a negative literal contributes $C_i(Y^{D/\ell_i})$. Append the $n$ guards and enumerate the resulting polynomials as $P_0,\ldots,P_{M-1}$. At every $D$th root their values lie in $\mathbb F_2$, and they vanish simultaneously exactly at the representatives of satisfying assignments.

Choose a prime $q$ different from the $\ell_i$ with $\operatorname{ord}_q(2)>M$. Define
\begin{align}
 N&=Dq,& B(X)&=X^N-1,\\
 E(X)&=\sum_{h=0}^{q-1}X^{Dh},&
 g(X)&=\sum_{j=0}^{M-1}X^{Dj}P_j(X^q),\\
 H(X)&=E(X)+(1+E(X))g(X),&
 T(X)&=H(X)^4+XB(X)^2.
\end{align}
The symbols $S$ and $N$ have different roles: $S$ is a witness count, whereas $N$ is the order of a roots-of-unity group.

\begin{theorem}\label{thm:ff-binary}
The polynomials above have polynomial-size sparse binary-exponent descriptions. The finite scheme
$Z=V(H,B)\subset\mathbb A^1_{\mathbb F_2}$ is finite \'etale, possibly empty, and is the reduced repeated-root locus of $T$. It has exactly $KS$ geometric points, all of multiplicative order $N$, where
\[
 K=(q-1)\prod_i\frac{\ell_i-1}{2},
 \qquad
 L=\operatorname{ord}_N(2)
   =\operatorname{lcm}\bigl(\operatorname{ord}_q(2),
                            \operatorname{ord}_{\ell_i}(2):1\leq i\leq n\bigr).
\]
All repeated roots of $T$ have multiplicity exactly two, all other roots are simple, and every closed point of $Z$ has degree $L$. In particular $L\mid K$.
\end{theorem}
\begin{proof}
For $\alpha\in\mu_N$, write $y=\alpha^q$ and $\beta=\alpha^D$. This identifies $\mu_N$ with $\mu_D\times\mu_q$. The polynomial $E$ equals one for $\beta=1$ and zero otherwise. Consequently $H(\alpha)=1$ when $\beta=1$. Otherwise
\[
 H(\alpha)=\sum_{j=0}^{M-1}P_j(y)\beta^j.
\]
The minimal polynomial of $\beta$ over $\mathbb F_2$ has degree $\operatorname{ord}_q(2)>M$. No nonzero polynomial of degree at most $M-1$ with coefficients in $\mathbb F_2$ can vanish there. Hence $H(\alpha)=0$ if and only if all the tests vanish. The balanced fibers give $KS$ roots. Their two coordinates have orders $D$ and $q$, so each root has order $N$.

Since $N$ is odd, $B$ is squarefree. Thus $\mathbb F_2[X]/(B,H)$ is a product of finite separable fields, proving the finite \'etale assertion. In characteristic two $T'=B^2$. At a common root $\alpha$ of $H$ and $B$, the summands $H^4$ and $XB^2$ have vanishing orders at least four and exactly two. Their sum therefore has order exactly two. A root of $T$ outside $B=0$ has nonzero derivative. This also proves
\[
 \gcd(T,T')=\gcd(H,B)^2.
\]
Frobenius acts on a primitive $N$th root by squaring, with orbit length $L$. Each $\operatorname{ord}_{\ell_i}(2)$ divides $(\ell_i-1)/2$, and the tag order divides $q-1$; hence their lcm divides $K$.

For completeness, all searches have deterministic polynomial bit cost. The prime number theorem in the fixed progression $1\pmod8$ bounds the first $n$ eligible primes by $O(n\log(n+2))$. A sieve with successively doubled bounds finds them. At most $M(M+1)/2$ distinct odd primes divide $\prod_{j=1}^M(2^j-1)$: the product has fewer than that many binary digits. Testing the first $M$ powers of two at a polynomial number of small prime candidates therefore finds an eligible $q$, with polynomial numerical size. Orders at these small primes are computed by repeated multiplication, and $L$ by gcd and lcm.

Let $\ell_{\max}=\max_i\ell_i$. Expanding an individual clause uses at most $(1+\ell_{\max})^3$ terms. The $q$-term selector increases the total number of terms only by another polynomial factor. Exponents in $g$ are less than $D(M+3q)$; in $H$ and $T$ they still have $O(\log D+\log(M+q))$ bits. Frobenius fourth powers merely multiply exponents by four. Combining coincident exponents is polynomial-time sorting and parity addition. Finally $N$, $K$, and $L$ all have $O(\log N)$ bits, though their numerical values can be exponential.
\end{proof}

The constant coefficient of $T$ is one, its only odd exponents are $1$ and $2N+1$, and $T(1)=1$. Indeed $E(0)=E(1)=1$ gives $H(0)=H(1)=1$, while every exponent of $H^4$ is divisible by four. These checks also distinguish the construction from an input consisting of a dense gcd: the expanded monic polynomial $\gcd(H,B)$ can have an exponentially different size.

\begin{corollary}\label{cor:ff-zeta}
For every positive integer $s$,
\[
 a_s:=\#Z(\mathbb F_{2^s})=KS\,\mathbf 1_{L\mid s},
 \qquad
 Z(Z,t)=(1-t^L)^{-KS/L}.
\]
For a prime $\lambda\neq2$, each $L$th root of unity occurs as a Frobenius eigenvalue on $H^0_{\mathrm{\acute et}}(Z_{\overline{\mathbb F}_2},\overline{\mathbb Q}_\lambda)$ with multiplicity $KS/L$. If $S=0$, $Z$ is empty, $H^0=0$, and its zeta function is one.
\end{corollary}
\begin{proof}
The scheme is the disjoint union of $KS/L$ spectra of $\mathbb F_{2^L}$. Each contributes $L$ embeddings into $\mathbb F_{2^s}$ when $L\mid s$, and none otherwise. Thus
\[
 \sum_{s\geq1}\frac{a_s}{s}t^s
  =\frac{KS}{L}\sum_{j\geq1}\frac{t^{Lj}}j
  =-\frac{KS}{L}\log(1-t^L).
\]
On geometric points Frobenius is a disjoint union of cycles of length $L$. The characteristic polynomial of one cycle is $U^L-1$, which splits with distinct roots in characteristic zero. This proves the eigenvalue multiplicities, including the case where $\lambda$ divides $L$.
\end{proof}

\subsection{What a weight estimate leaves undetermined}
Theorem~1(b) of \emph{A short proof of Deligne's bound for Frobenius eigenvalues, after Alp\"oge} states the usual purity conclusion for smooth proper schemes; its Proposition~4, p.~5, expresses a remainder with a single sign \cite{ff-sign}. Write $\mathbb G_m=\operatorname{Spec}\mathbb F_2[u,u^{-1}]$ for the multiplicative group. To test that formula in positive dimension, put
\[
 U=\mathbb G_m\times Z,\qquad Y=\{1\}\times Z,
\]
use the constant $\mathbb Q_\lambda$-sheaf. If $Z\neq\varnothing$, $U$ is smooth affine of pure dimension one and $Y\subset U$ has codimension one. On every geometric component, restriction $H^0(\mathbb G_m)\to H^0(\{1\})$ is an isomorphism. By duality the corresponding Gysin map
$H_c^0(Y)(-1)\to H_c^2(U)$ is an isomorphism. Thus the Gysin kernel in that proposition is zero. Writing
\[
 M_s=(2^s-1)a_s,\qquad N_s=a_s,
\]
its remainder is exactly
\[
 2^sN_s-M_s=a_s
 =\operatorname{Tr}\bigl(F^s\mid H_c^1(U,\mathbb Q_\lambda)\bigr).
\]
Indeed $H_c^1(\mathbb G_m,\mathbb Q_\lambda)=\mathbb Q_\lambda$ with trivial Frobenius, and K\"unneth identifies $H_c^1(U)$ with the permutation representation on $Z$. When $Z$ is empty, all terms vanish directly and no pure-dimensional hypothesis is needed. The dimension-zero argument in Section~8, p.~8, of the same paper gives the identical finite permutation representation. Every eigenvalue has absolute value one; the unknown multiplicity $KS/L$ of each cycle eigenvalue remains. This is a calculation inside the weight formalism, not a failure of that formalism: the absolute-value estimate contains no procedure for recovering this multiplicity from sparse equations.

\subsection{An invertible normalization for every fixed prime}
Fix a prime $p$. Choose distinct primes $\ell_1,\ldots,\ell_n,q$ all congruent to $-1\pmod{4p}$ and require $\operatorname{ord}_q(p)>M$. Put
\begin{gather*}
 D=\prod_i\ell_i,\quad N=Dq,\quad B(X)=X^N-1,\\
 K_* =\frac{q-1}{2}\prod_i\frac{\ell_i-1}{2},\quad
 L=\operatorname{ord}_N(p).
\end{gather*}
For one of these primes $a$, let $R_a$ be its nonzero square residues and set
\[
 C_a(U)=-\sum_{v\in R_a}U^v,
 \qquad
 \Gamma_a(U)=a^{-1}\sum_{v=0}^{a-1}U^v
 \quad\text{in }\mathbb F_p[U].
\]
These are again Boolean bits and guards on $\mu_a$. Here is a direct verification, including $p=2$. For a primitive $a$th root $u$, let $\eta_0,\eta_1$ be the two quadratic periods. Their sum is $-1$. Since $a\equiv3\pmod4$, negation exchanges the two classes. In $\eta_0\eta_1$, zero occurs as an exponent $(a-1)/2$ times. Multiplication by any nonzero scalar preserves the two classes or exchanges them, so every nonzero exponent occurs equally often, namely $(a-3)/4$ times. Hence
\[
 \eta_0\eta_1=\frac{a-1}{2}-\frac{a-3}{4}
             =\frac{a+1}{4}=0\quad\text{in }\mathbb F_p.
\]
All quotients here are integers before reduction. Therefore the periods are $0,-1$. Nonsquare powers exchange them, so $C_a$ takes each of $0,1$ on precisely $(a-1)/2$ primitive roots. At $1$, $C_a(1)=-(a-1)/2=1$ in $\mathbb F_p$, while $\Gamma_a$ is the indicator of $1$.

Build the Boolean violation polynomials using $1-C_{\ell_i}$ for positive literals and $C_{\ell_i}$ for negative literals, followed by the guards; again call them $P_j$. Define
\begin{align}
 g(X)&=\sum_{j=0}^{M-1}X^{Dj}P_j(X^q),\\
 W(X)&=(1-\Gamma_q(X^D))C_q(X^D),\\
 A(X)&=\bigl(1-W(X)+W(X)g(X)\bigr)\bmod (X^N-1),
 \qquad \deg A<N.
\end{align}
At a passing tag $W=1$, the same degree argument forces all $P_j$ to vanish; at every other tag $A=1$. Thus precisely half of the primitive tag roots are allowed.

\begin{theorem}\label{thm:ff-unit}
The pair $(A,B)$ has a deterministic polynomial-size sparse description for fixed $p$. Its common roots are exactly $K_*S$ distinct primitive $N$th roots, and $A(1)=1$. Moreover $L\mid K_*$ and $p\nmid K_*$. The homogeneous equations
\[
 V^NA(X/V)=0,\qquad X^N-V^N=0
 \quad\text{in }\mathbb P^1_{\mathbb F_p}
\]
define a finite \'etale scheme $Z$ with $(K_*/L)S$ closed points, each of degree $L$. It is empty when $S=0$ and geometrically disconnected whenever $K_*S>1$.
\end{theorem}
\begin{proof}
The preceding balanced-fiber calculation gives the root count and primitivity. Frobenius preserves each Boolean fiber, since $C_a$ has coefficients in $\mathbb F_p$. Its orbits among primitive $a$th roots have size $\operatorname{ord}_a(p)$, so that order divides $(a-1)/2$. Taking their lcm proves $L\mid K_*$. Every factor $(a-1)/2$ is $-1\pmod p$, hence $K_*$ is a unit modulo $p$; so are $L$ and $K_*/L$.

The second projective equation excludes infinity. On $V\neq0$ the coordinate algebra is $\mathbb F_p[X]/(A,B)$; since $B'=NX^{N-1}$ has no common root with $B$, this algebra is finite \'etale. All its geometric points have Frobenius orbit length $L$, proving the residue-field assertion.

The prime search is polynomial for fixed $p$: the progression has fixed modulus $4p$, and at most $M(M+1)\log_2(p)/2$ distinct primes divide $\prod_{j=1}^M(p^j-1)$. A polynomial number of progression candidates therefore suffices, after exclusion of the variable primes. Individual clauses and the tag product have polynomially many terms, with the latter contributing $O(q^2)$ before cancellation. Remainder modulo $X^N-1$ is computed simply by reducing each binary exponent modulo $N$ and combining equal terms. Homogenization records exponent pairs $(e,N-e)$ and preserves polynomial bit size.
\end{proof}

The polynomial $A^{p^2}+XB^p$ has precisely these repeated roots, each of multiplicity $p$: its derivative is $B^p$, and the two summands have orders at least $p^2$ and exactly $p$ at each common root. The binary construction of Theorem~\ref{thm:ff-binary} uses the full nonidentity tag and has its own normalization $K$; it should not be confused with $K_*$.

\subsection{Genus, geometric $p$-rank, and coherent trace}
The sparse affine equations
\begin{equation}\label{eq:ff-AS}
 B(X)W=1,\qquad Y^p-Y=A(X)W
\end{equation}
define a smooth affine curve: the first equation localizes the affine line, and the second has derivative $-1$ in $Y$. For the original binary construction, take $A=H\bmod B$, with $\deg A<N$, in these same equations. Let $C$ be the smooth projective completion. This presentation is an affine open together with a function-field convention; it is not a polynomial-size list of ordinary projective equations for the completion.

\begin{theorem}\label{thm:ff-genus}
For the normalization of Theorem~\ref{thm:ff-unit}, $C$ is geometrically integral and its genus and geometric $p$-rank are both
\[
 g(C)=f_p(C)=(p-1)(N-1-K_*S).
\]
For the original binary construction the same formula holds with $K$ in place of $K_*$. In the unit normalization either invariant, even modulo $p$, determines $S\bmod p$ by
\[
 S\equiv K_*^{-1}\bigl(g(C)+N-1\bigr)\pmod p.
\]
\end{theorem}
\begin{proof}
The rational function $A/B$ has a simple pole at precisely those roots of $B$ where $A$ is nonzero. There are $b=N-K_*S$ such poles. The other roots are removable, because $B$ is squarefree. Since $\deg A<N$, there is no pole at infinity; since $A(1)=1$, at least the pole at $1$ remains. A function $h^p-h$ over $\overline{\mathbb F}_p(X)$ has all pole orders divisible by $p$. Thus $A/B$ cannot be of that form. The Artin--Schreier polynomial is irreducible: its splitting field is obtained by adjoining any one root, and its Galois group is a subgroup of the additive group $\mathbb F_p$; failure of irreducibility would make that group trivial and give a root in the rational-function field. Hence the extension has degree $p$ even geometrically.

Every pole has lower ramification jump one, and different exponent $2(p-1)$. Riemann--Hurwitz gives
$2g-2=-2p+2b(p-1)$, so $g=(p-1)(b-1)$. The Deuring--Shafarevich formula gives $f_p=(p-1)(b-1)$. These are the specializations $e_j=2$ of Pries and Zhu (2012, Lemma~2.6, p.~712) \cite{ff-pries-zhu}. Reduction modulo $p$ gives $g\equiv K_*S-N+1$, proving the decoding formula.
\end{proof}

For $\lambda\neq p$, geometric $H^1_{\mathrm{\acute et}}(C,\mathbb Q_\lambda)$ has dimension $2g$, and all its Frobenius eigenvalues have absolute value $\sqrt p$. The degree of the weight-one zeta factor therefore determines $S$ exactly, just as the genus does. At $p=2$, recovery of genus parity from the Betti number requires that number modulo four, since the Betti number is $2g$. The magnitude statement alone gives no value of $g$.

For the finite projective scheme there is a sharper contrast. With $r=(K_*/L)S$,
\[
 H^0(Z,\mathcal O_Z)\simeq\prod_{j=1}^r\mathbb F_{p^L}.
\]
The normal basis theorem gives a basis of each factor cyclically permuted by the $p$-power map $F$. Hence, for every $a\geq1$,
\begin{equation}\label{eq:ff-trace}
 \operatorname{Tr}(F^a)=K_*S\,\mathbf 1_{L\mid a}\quad\text{in }\mathbb F_p.
\end{equation}
In particular ordinary $F$ has trace zero because $L>M\geq2$, whereas $F^L$ is the identity and has trace $K_*S$. The binary integer $L$ and the two sparse equations specify this iterated-trace problem in polynomial bit length. Its coherent dimension $K_*S$ may be exponential.

As a control, homogenize the Boolean equations $x_i(x_i-1)=0$ and clause-violation products to obtain
\[
 X_i(X_i-V)=0,\qquad
 \prod_{\ell\in C}(V-\ell^h)=0
 \quad\text{in }\mathbb P^n.
\]
For a negative literal $1-x_i$, its homogeneous literal is $V-X_i$; for a positive literal it is $X_i$. Empty clauses give $1=0$ and an empty scheme. At $V=0$, the Boolean equations force all $X_i=0$, so infinity is absent scheme-theoretically after projectivization. On $V=1$ the Boolean coordinate ring splits as a product indexed by bit words; the clause equations delete exactly the nonsatisfying factors. The result is the disjoint union of $S$ rational points, with $F$ the identity on an $S$-dimensional coherent space and trace $S\bmod p$. Thus increasing ambient dimension gives the uniterated trace immediately, but it does not give a small coherent space. Neither this construction nor the curve construction replaces the coherent line of the smooth circuit hypersurface.

\subsection{Numerical examples and a varying-field curve}
For $p=2$, take $\ell_1=7$ and $q=23$. The powers of two modulo $7$ are $1,2,4$, and $2^{11}\equiv1\pmod{23}$ with no smaller positive exponent. Thus
\[
 N=161,\qquad K_*=3\cdot11=33,\qquad L=33.
\]
The square residues modulo $7$ are $1,2,4$; those modulo $23$ are
$1,2,3,4,6,8,9,12,13,16,18$.
For the positive unit clause, the two tests are $1-C_7$ and $\Gamma_7$, so
\[
 g(X)=1-C_7(X^{23})+X^7\Gamma_7(X^{23}),\qquad
 A=(1-W+Wg)\bmod(X^{161}-1).
\]
Exactly $33$ primitive roots pass. For the contradictory, unit, and tautological one-variable formulas, $S=0,1,2$ respectively; their genera are $160,127,94$. Each satisfying assignment contributes one closed point of degree $33$ and the genus parity equals $S\bmod2$.

For $p=3$, take $\ell_1=11$, $q=23$. The successive powers of three modulo $11$ have order five, and $3^{11}\equiv1\pmod{23}$ with order eleven. Thus $N=253$ and $K_*=L=55$. The three genera are $504,394,284$, congruent to $0,1,2$ modulo three. Their zeta functions as finite schemes are $(1-t^{55})^{-S}$, and $\operatorname{Tr}(F^{55})=55S\equiv S\pmod3$. Ordinary trace is zero in all three cases. The quadratic-period proof, rather than a root list in an exponentially large field, verifies these counts.

There is also a smooth projective curve encoding the exact count over a varying binary field. We record the construction because it shows why the whole coherent Frobenius operator need not recover the next counting digit. Let $b$ be the least power of two with $b\geq n$, put $q=2^b$, and construct $K=\mathbb F_q$ by successive quadratic extensions. At each step, the absolute trace is a nonzero linear functional (its polynomial has degree smaller than the field size), so one can find a trace-one element by testing a basis and adjoin a root of $U^2+U+a$. This gives deterministic polynomial-time arithmetic and, if desired, a power-basis defining polynomial.

Take trace-dual bases $e_i,\beta_i$ and set
\[
 L_i(X)=\operatorname{Tr}_{K/\mathbb F_2}(\beta_iX)
       =\sum_{j=0}^{b-1}(\beta_iX)^{2^j}.
\]
If $x=\sum a_ie_i$, then $L_i(x)=a_i$. Let $s_\Phi$ be the Boolean satisfaction polynomial, obtained by multiplying the clause-success polynomials. Put
\[
 P(X)=1+s_\Phi(L_1(X),\ldots,L_n(X))
                  \prod_{i=n+1}^b(1+L_i(X)).
\]
It is zero on exactly $S$ elements of $K$. With $\operatorname{Tr}(\lambda)=1$, choose the odd integer
$D_0=(3m+b+1)q+1$ and set
\[
 f(X)=\lambda P(X)+(X^q+X)X^{D_0-q}.
\]
Its degree is exactly $D_0$, it is monic, and its circuit and binary degree have polynomial bit length. On $K$, $\operatorname{Tr}(f(x))=P(x)$, so the equation $y^2+y=f(x)$ has exactly $2S$ affine $K$-points.

Write $D_0=2g+1$, $w=g+1$. Its completion
\[
 C_\Phi:\quad Y^2+YZ^w=Z^{D_0+1}f(X/Z)
       \subset\mathbb P_K(1,w,1)
\]
is smooth and geometrically integral and has exactly one point at infinity. Indeed the charts $Z\neq0$ and $X\neq0$ cover it. On the first the $y$ derivative is one; on the second, with $z=Z/X$ and $v=Y/X^w$, the equation is
$v^2+vz^w=z^{D_0+1}f(1/z)$, whose right side begins with $z$. The unique point $(z,v)=(0,0)$ is therefore smooth. Irreducibility follows from the odd pole order of $f$ at infinity. Consequently $\#C_\Phi(K)=1+2S$.

The two affine charts compute $H^1(\mathcal O)$ with basis
$e_j=[x^{-j}y]$, $1\leq j\leq g$. Constant Laurent polynomials are coboundaries, while the $y$-coefficients are quotiented by
$K[x]+x^{-w}K[x^{-1}]$. Hence absolute Frobenius and $q$-Frobenius are respectively
\[
 F(e_j)=\begin{cases}e_{2j}&2j\leq g,\\0&2j>g,\end{cases}
 \qquad
 F^b(e_j)=\begin{cases}e_{qj}&qj\leq g,\\0&qj>g.\end{cases}
\]
They are nilpotent and depend only on the degree and field. The same genus and field can therefore accompany different values of $S$. The coherent zeta congruence of SGA~7~II, Expos\'e~XXII, Theorem~3.1 \cite{ff-sga7}, is always $(1-t)^{-1}$ modulo two, but the integer \'etale trace is $q-2S$, so the next binary digit recovers the parity of $S$. The degree, genus, and ordinary projective Veronese embedding may all be exponential. This construction concerns circuit equations in weighted projective space, and it is not a bounded-Betti-number solution of the small-zeta question.

\begin{proposition}\label{prop:ff-coherent-limit}
On the varying-field curve family just constructed, the field, genus, and complete based coherent Frobenius operator do not determine $S$ modulo any integer $m>1$. This remains true when the operator is represented by the displayed index rule with binary parameters $g,q$ rather than by an expanded matrix.
\end{proposition}
\begin{proof}
Use two declared variables and one clause. The formulas $x_1$ and $x_1\vee x_2$ have respectively two and three satisfying assignments. The construction has $b=2$, $q=4$, $D_0=(3+2+1)4+1=25$, and $g=12$ for both formulas. Their field presentation, genus, and Frobenius rules $e_j\mapsto e_{2j}$ and $e_j\mapsto e_{4j}$ therefore agree. Their counts differ by one, which is nonzero modulo every $m>1$. Consequently any function of these identical invariant data has the same output on the two curves and cannot return both correct residues.
\end{proof}

This limitation is specific to the stated invariant and presentation, and is unconditional. It does not say that those two curves have the same full zeta function: their ground-field counts are five and seven. In particular it does not convert an invariant that is easy to describe into an algorithm for the integer \'etale trace.
